\section{Verification and formal boundary}\label{sec:verification}

No geometric classification theorem in this paper is claimed to be
formalized in Lean.  The formal layer checks coordinate identities and
conditional synthesis implications with the component geometry,
rational-point existence, group actions, and certificate semantics
kept as explicit inputs.

The conceptual theorems above are accompanied by deterministic evidence
bundles.  Table~\ref{tab:evidence} records which claims use exhaustive
computation.  Each public label resolves, through the supplement
schema, to its generator, certificate, replay where stated, and
SHA-256 manifest.  Table~\ref{tab:dependency-map} summarizes the corresponding
mathematical mechanisms and closing gates.

\small
\begin{longtable}{>{\raggedright\arraybackslash}p{0.19\textwidth}
                  >{\raggedright\arraybackslash}p{0.25\textwidth}
                  >{\raggedright\arraybackslash}p{0.47\textwidth}}
\caption{Public verification map.  Exact filenames, commands, hashes,
domains, and stop conditions belong to the electronic supplement.}
\label{tab:evidence}\\
\toprule
Public label & Verified domain & Replay and trust boundary\\
\midrule
\endfirsthead
\toprule
Public label & Verified domain & Replay and trust boundary\\
\midrule
\endhead
\bottomrule
\endfoot
Certificate R5 & nineteen audited field records; seven-field finite
bridge, with \(q=16\) covered as a regression check rather than a
dependency & independent pointwise and orbitwise replay;
canonical representatives and exhaustion records\\
Certificate R6 & direct scan through $q=16$; finite bridge below $29$ &
independent syndrome replay and separate radius gate\\
Certificate R6-NF & small exceptional normal forms & deterministic
normal-form reconstruction; not a second implementation\\
Certificate R7 & $q=7,8,9,11$ and finite bridge below $37$ &
primary quotient enumeration and representative/orbit replay;
split-free and radius flags remain separate\\
Certificate R7 direct locus & all fourteen R7 certificate fields &
candidate-domain and marker-intersection reconstruction, orbit and
Frobenius replay; shares the direct-locus engine, R5 field layer, and
R6 pointed theorem from \(q\geq16\), and is not a second field
implementation\\
Certificate SC & terminal factorization and coherent-Fano identities &
Python and Singular elimination replays; density and finite-component
selection are separately kernel checked\\
Certificate R8 & threshold, nucleus, witness, and budget data &
generator plus separately written replay; no ambient syndrome census\\
Certificate R9 & residual-quadratic and characteristic-seven data &
independent algebra replay and exact \(q=49\) carrier bridge\\
Certificate R10 & threshold and persistent orbit arithmetic &
separate prime-power and cyclic-orbit replay\\
Certificate Lucas M9 & rank-two twists; full \(q=16,32\) carrier;
\(q=64\) invariant block & independent action and Hankel replay; the
\(q=64\) complement and the uniform genus-one argument are mathematical\\
\end{longtable}
\normalsize

Table~\ref{tab:r6-exceptional} is a deterministic projection of the
hash-pinned public R6 and R6-NF records.  The top-level verifier
regenerates it and checks byte-for-byte equality with the included
source.

The supplement's \texttt{CERTIFICATE-SCHEMA.md} defines the public labels and
replay classes; its reproduction guide and two manifests record commands,
locks, hashes, and byte counts. The entry point \texttt{supplement/verify.py}
checks those records, the classification projections, and the local PDF;
\texttt{--replay} runs the paper-local Python replays.

The exact Lean declarations, hypotheses, axiom audit, and replay command are
listed in \texttt{LEAN-STATEMENTS.md}. They check the coordinate,
budget, density, closure, component-selection, and conditional-synthesis
interfaces. Concrete component geometry, fixed-level carrier arithmetic,
classical radius inputs, and certificate semantics remain explicit inputs.
The formal layer also checks the field-eight record but draws no quantum
consequence from it, because the required one-column MDS extension does not
exist at radius \(r-1\) (Remark~\ref{rem:q8-no-extension}).

\subsection{Data, code, and disclosure}\label{sec:provenance}

The electronic supplement accompanying this manuscript contains the
source, canonical classification representatives and stabilizers,
certificate schemas, generators, independent replays where available,
expected hashes, toolchain locks, and a top-level verification command.
This appendix assigns each computational or formal claim its exact trust
route.

The repository also contains Projective Reed--Solomon Toolkit under the
\texttt{software/} directory. Its command-line interface separates structural
computation from theorem-gated classification.  For every $r\geq5$ and
$q\geq r$, \texttt{canonicalize} returns an exact semilinear canonical form
without a coding verdict, while \texttt{distance} and \texttt{decode} return
exact nearest-error data and a replayable locator certificate within the
stated candidate budget.  The \texttt{classify} command applies the frozen
R5--R10 theorem registry and returns \texttt{DEEP}, \texttt{NOT\_DEEP},
\texttt{UNRESOLVED}, or \texttt{UNSUPPORTED}; \texttt{verify} independently
replays emitted locator and positive deep certificates.

Locator candidates are streamed, so the candidate budget bounds search time
rather than allocated locator storage. The broader canonicalization and
decoding domains do not promote an R11-or-higher syndrome to a deep hole
without an independently proved covering radius and family exhaustion.  For
the one enabled higher-redundancy family, even $q$ and $r=q-1$, the verifier
replays the $e_{r-2}$ terminal-locator obstruction and imports the
dimension-two radius theorem of \cite[Thm.~17(1)]{WuDingChen2023}.  This is a
certified family, not a complete R11-or-higher classification.
